% !TEX root = ../../vzor.tex
%% Ukázka: číslované rovnice, align, řecká písmena, vektory
%% Vkládá ji soubor vzor/kapitoly/teorie.tex; ve vzor.pdf je pod nadpisem uvedeno „Zdroj ukázky".

\section{Matematické výrazy}
\zdrojukazky{vzor/ukazky/matematika.tex}

Číslovaná rovnice se vytvoří prostředím \texttt{equation}; odkazujeme na ni příkazem \verb|\eqref{}|.
Pythagorova věta~\eqref{eq:pythagoras}:
\begin{equation}
	a^2 + b^2 = c^2
	\label{eq:pythagoras}
\end{equation}

Víceřádkové odvození se zarovnaným rovnítkem (prostředí \texttt{align}):
\begin{align}
	E         &= mc^2 \\
	\vec{F}   &= m\vec{a} \\
	W         &= \vec{F} \cdot \vec{s}
\end{align}

Příklady zápisů v textu: řecká písmena $\alpha, \beta, \omega, \Delta, \lambda$,
integrál $\int_a^b f(x)\,dx$, parciální derivace $\frac{\partial f}{\partial x}$,
vektory $\vec{F}$, normálový vektor $\hat{n}$.
